% 第一章 绪论、文献综述与研究方法

\section{绪论、文献综述与研究方法}

\subsection{引言}

股票市场收益率与宏观经济变量之间的关系，始终是金融研究中的重要议题。经典研究已经表明，股票预期收益并非固定不变，而会随着商业周期、风险溢价与整体经济状态的变化而调整。\cite{fama1989business}发现，股票与长期债券的预期收益呈现出清晰的商业周期特征，经济环境偏弱时预期收益较高，经济状况偏强时预期收益相对较低，这说明宏观状态能够系统性地进入资产定价过程。随后，大量研究继续围绕收益预测问题展开讨论。\cite{campbell2008predicting}认为，在预测回归中加入适度的经济约束后，部分变量仍能在样本外提供有限但具有经济意义的预测增益；\cite{welch2008comprehensive}则指出，许多变量虽然在样本内表现突出，却难以持续稳定地优于历史均值模型。这些研究共同提示，宏观变量更适合作为收益率条件预期的构造基础。

这一问题在中国资本市场中具有更强的现实意义。A股市场长期表现出较高的政策敏感性、情绪波动性和资金驱动特征，宏观变量、政策预期、汇率变化以及市场风险情绪之间往往存在更为紧密的互动。随着北向资金持续进入、期权市场逐步扩容以及高频交易信息不断丰富，市场价格形成机制已经同时包含宏观基本面、跨境资金流、投资者情绪和流动性摩擦等多重因素。近年的中国研究也表明，宏观经济不确定性在A股市场中存在显著定价效应，北向资金会改变市场流动性和价格表现，投资者情绪及信息环境也会进一步影响收益形成机制\cite{ge2024macroeconomic,yang2023northbound,he2024timevarying}。在这一背景下，研究宏观信息如何形成收益预期、市场为何偏离这一预期，具有明显的理论价值和现实意义。

基于上述认识，本文将研究重点放在两个相互衔接的环节。第一阶段利用月度宏观经济变量构建沪深300指数未来5日与60日累计收益率的条件预期；第二阶段以实际收益率与条件预期收益率之差构造异常收益及其绝对值，并进一步检验市场情绪、风险感知、资金流向、杠杆交易、流动性和交易状态等因素对收益偏离的解释作用。这样的处理能够将"基本面信息进入收益率"的过程与"市场围绕基本面发生偏离"的过程区分开来，也有助于更清楚地刻画短线与长线收益波动的形成机制。

\subsection{研究现状与文献综述}

\subsubsection{宏观经济变量与股票收益率预测研究}

宏观经济变量能否预测股票收益率，是资产定价与经验金融研究中的经典问题。\cite{fama1989business}较早将股票收益与商业周期联系起来，指出预期收益具有明显的宏观状态依赖性。在此基础上，\cite{campbell2008predicting}从样本外预测角度重新讨论了收益预测问题，认为若在模型中施加合理的经济约束，部分预测回归在样本外仍可取得一定改进。\cite{welch2008comprehensive}则系统比较了权益溢价预测变量的经验表现，发现许多变量在样本内看似有效，但在样本外往往难以持续优于历史均值。由此形成的基本共识是，宏观变量对股票收益率的作用更多体现为条件预期的变化，而非一种简单、固定、线性的决定关系。

国内研究同样表明，宏观变量在中国股票市场中具有重要作用，但这种作用往往带有更强的状态依赖性和阶段差异。\cite{ge2024macroeconomic}基于中国A股市场数据构建宏观经济不确定性指标后发现，宏观不确定性暴露存在显著溢价效应，而且这一效应会随公司特征和市场环境变化而发生异质性调整。这些研究说明，在中国市场中讨论宏观变量与股票收益率的关系时，单一平均效应往往不足以捕捉真实机制，将宏观信息置于"条件预期收益率"的框架中进行处理，更符合当前经验研究的发展方向。

\subsubsection{混频数据模型与 MIDAS 方法研究}

宏观变量预测金融市场收益率时面临的直接难题，在于数据频率不一致。宏观指标通常以月度或季度发布，而股票收益率则以日度持续变动。如果直接将低频变量机械复制到高频样本中，既无法充分体现宏观信息的动态传导，也会扭曲信息到达顺序。\cite{ghysels2004midas,ghysels2006midas}提出的 MIDAS 模型，为不同频率数据的联合建模提供了系统方法。该模型通过参数化权重函数，将低频变量的多个滞后期压缩到高频回归方程中，从而在保留时序结构信息的同时避免参数维数过快膨胀。随后，\cite{ghysels2007midas}对 MIDAS 回归的估计和应用进行了进一步拓展，说明其不仅适用于波动率预测，同样适用于一般意义上的收益率和宏观变量预测。\cite{schorfheide2015realtime}则从实时预测角度比较了 mixed-frequency VAR 与 MIDAS 回归，指出混频框架在处理宏观信息逐步披露时具有明显优势。

在中国语境下，混频方法也已经获得较为稳固的应用基础。\cite{zheng2013midas}将 MIDAS 模型用于中国 GDP 的短期预测，发现基于金融变量的 MIDAS 模型在实时数据背景下具有较好的预测表现，并且多元扩展模型在若干设定下优于基准模型。这一结果说明，MIDAS 方法在中国宏观金融数据环境中具有较好的可操作性。本文主要从两个方面使用这一方法：其一，月度宏观变量能够与日度收益窗口直接对接；其二，过去若干期宏观信息的影响强度可以通过权重函数灵活刻画，从而更贴合宏观冲击对市场收益率并非一次性释放的现实特征。

\subsubsection{投资者情绪、风险感知与收益偏离研究}

在经典资产定价框架之外，行为金融研究强调，投资者情绪会改变资产定价并推动价格偏离基本面。\cite{baker2006investor}的研究表明，情绪波动对那些估值主观性强、套利难度较高的资产影响更为显著，这意味着市场价格围绕基本面发生偏离，也与投资者风险偏好和情绪波动直接相关。

这一研究路径在中国市场得到了进一步扩展。\cite{wang2004investor}较早指出，投资者情绪会显著影响我国股市收益与收益波动。\cite{yin2019highfrequency}利用高频情绪指标发现，投资者高频情绪对股票日内收益率具有明显预测作用。\cite{mei2019investor}基于移动互联网文本构造情绪指标，发现乐观情绪越高，下一期股票收益越高，而且这种作用在信息环境较差、流动性较弱的股票上更为突出。相关研究表明，在中国股票市场中，情绪和风险感知会直接影响收益率的实现路径。对于本文的研究设计而言，情绪变量更适合放在第二阶段，用于解释市场为何围绕宏观条件预期发生偏离。

\subsubsection{资金流向、流动性与交易状态研究}

除情绪因素外，资金流向和市场流动性同样会显著影响收益率表现。\cite{amihud2002illiquidity}从时间序列和横截面两个层面说明，预期市场非流动性上升会提高投资者要求的回报补偿，股票收益中包含明确的流动性溢价。这一结果在后续研究中得到多次验证，也使流动性成为理解收益偏离和市场异动的重要视角。

在中国市场中，北向资金的作用受到越来越多关注。\cite{yang2023northbound}基于高频数据研究发现，北向资金整体上能够提升股票流动性，其作用机制主要体现在信息机制和竞争机制，但在特定情境下也可能通过价格冲击和筹码约束影响市场质量。\cite{he2024timevarying}进一步指出，北向资金与中国股市表现之间存在时变互动关系，其对收益的影响具有明显短期性，并在极端波动时期呈现出一定的稳定市场特征。这些研究提示，资金流向、杠杆交易与流动性状态不仅影响收益率本身，也会影响市场围绕宏观基本面形成预期偏离的幅度和持续性。

\subsubsection{文献评述}

现有研究已经形成较为完整的理论与方法基础。宏观收益预测文献关注经济状态如何进入条件预期，MIDAS 研究解决了低频宏观变量与高频收益率之间的频率衔接问题，情绪、资金与流动性研究则揭示了市场围绕基本面发生偏离的多种机制。现有成果仍留下几处值得继续推进的空间。首先，许多研究聚焦月度或季度收益率，对未来5日与60日累计收益率这种"短线—长线并置"的预测窗口关注不足。其次，宏观收益预测研究与市场偏离解释研究往往分开展开，前者着重于预测能力，后者偏重于行为和微观结构因素，较少在同一框架中先构造宏观条件预期，再单独检验非基本面因素如何放大偏离。最后，针对中国股票市场的研究虽然已经涉及宏观不确定性、投资者情绪、北向资金和流动性等关键变量，但把这些因素纳入统一的两阶段实证框架中进行比较分析的研究仍相对有限。本文的研究路径正是在上述文献基础上的进一步延伸。

\subsection{研究方法与技术路线}

本文采用两阶段实证框架。第一阶段围绕"宏观变量如何形成收益预期"展开，第二阶段关注"市场为何偏离这一预期"。

在第一阶段，本文选取 CPI、PPI、M2增速、经济政策不确定性指数和美元兑人民币汇率等月度宏观变量作为核心信息集，利用 MIDAS 混频回归模型预测沪深300指数未来5日和60日累计收益率。考虑到宏观变量对未来收益率的影响可能具有不同的滞后分布，本文使用参数化权重函数对过去12个月的宏观信息进行加权，并分别构建单变量 MIDAS 模型。在此基础上，再根据样本内拟合效果、样本外递推表现、VIF 检验结果以及变量的经济含义，筛选出最多4个变量进入精简多变量模型，形成未来5日和60日累计收益率的条件预期值。样本外预测采用扩展窗口法，以前60\%的样本作为初始估计区间，对后40\%的样本进行动态递推，并使用 RMSE、MAE 和样本外 $R^2$ 等指标评价模型表现\cite{ghysels2004midas,zheng2013midas}。

在第二阶段，本文根据第一阶段生成的条件预期收益率构造异常收益和异常收益绝对值。未来 $h$ 日异常收益定义为
\begin{equation}
AR_t^{(h)}=R_t^{(h)}-\widehat{R}_t^{(h)},\qquad h\in\{5,60\}
\label{eq:abnormal_return}
\end{equation}
同时定义异常收益绝对值
\begin{equation}
AbsAR_t^{(h)}=\left|AR_t^{(h)}\right|
\label{eq:abs_abnormal_return}
\end{equation}
其中，$AR_t^{(h)}$ 反映偏离方向，$AbsAR_t^{(h)}$ 刻画偏离强度。考虑到偏离强度更能反映市场异动和风险放大过程，本文在主结果中以 $AbsAR_t^{(5)}$ 和 $AbsAR_t^{(60)}$ 作为核心因变量，在补充分析中再讨论带方向的异常收益。第二阶段解释变量按经济含义划分为三组：情绪与风险感知变量，包括情绪标准分和隐含波动率指数；资金与杠杆交易变量，包括北向资金净流入和融资融券余额；流动性与交易状态变量，包括 Amihud 非流动性指标、20日动量和日内振幅。回归设计采用分组嵌套方式，先检验情绪与风险感知变量，再逐步加入资金与杠杆交易变量、流动性与交易状态变量，并在完整嵌套模型之后通过 LASSO 辅助筛选非零系数变量。新增变量组的边际解释力通过增量联合检验与调整后 $R^2$ 的变化共同判断。由于未来5日和60日收益窗口都属于重叠窗口，第二阶段回归统一采用 Newey--West HAC 稳健标准误进行修正\cite{he2024timevarying,amihud2002illiquidity}。

整个研究的技术路线可以概括为五个环节。第一步，基于沪深300指数收盘价构造未来5日和60日累计收益率，并将月度宏观变量按照"上一个完整月可得信息"原则与日度样本对齐。第二步，将样本划分为训练集与测试集，在训练集上完成缩尾阈值计算、标准化参数估计、VIF 检验和单变量 MIDAS 估计，并将相应处理规则外推到测试集。第三步，利用单变量与精简多变量 MIDAS 模型形成宏观条件预期收益率。第四步，在此基础上构造异常收益及其绝对值，形成市场偏离指标。第五步，采用分组嵌套回归、增量联合检验、LASSO筛选和稳健性分析识别情绪、资金、流动性和交易状态对偏离的解释作用。这一技术路线将宏观信息、收益预期和市场异动纳入同一分析框架，研究目标与方法设定之间的对应关系也更清晰。

% 图1-1：研究技术路线图
\begin{figure}[htbp]
\centering
\begin{tikzpicture}[
    node distance=0.8cm and 1.2cm,
    box/.style={rectangle, rounded corners, draw=blue!60, fill=blue!5, very thick, minimum width=3cm, minimum height=0.8cm, align=center, font=\small},
    arrow/.style={-{Stealth[length=2mm]}, thick, blue!60}
]

% 第一行
\node[box] (macro) {宏观变量输入\\CPI、PPI、M2、EPU、汇率};

% 第二行
\node[box, below=of macro] (midas) {MIDAS混频回归\\条件预期收益};

% 第三行
\node[box, below=of midas] (ar) {异常收益构造\\$AR_t^{(h)}$、$AbsAR_t^{(h)}$};

% 第四行
\node[box, below=of ar] (reg) {分组嵌套回归\\情绪/资金/流动性};

% 第五行
\node[box, below=of reg] (robust) {稳健性检验\\权重/滞后/子样本};

% 箭头
\draw[arrow] (macro) -- (midas);
\draw[arrow] (midas) -- (ar);
\draw[arrow] (ar) -- (reg);
\draw[arrow] (reg) -- (robust);

\end{tikzpicture}
\caption{研究技术路线图}
\label{fig:research_framework}
\end{figure}

\subsection{研究思路与预期贡献}

本文以宏观变量构造收益率条件预期为起点，再将市场状态变量纳入第二阶段解释框架，分别考察预期形成与预期偏离两个环节。宏观变量在模型中承担基本面信息集的作用，情绪、资金和流动性变量则用于识别市场围绕基本面出现偏离的机制。与单一方程同时放入全部变量相比，两阶段设计更便于区分变量功能，也更适合比较短线与长线收益窗口下作用路径的差异。

本文的可能贡献主要体现在三个方面。其一，在研究对象上，将未来5日和60日累计收益率同时纳入分析，在同一实证框架下比较宏观信息对短线与长线收益的条件预测能力。其二，在方法设计上，将 MIDAS 混频收益预测与异常收益偏离解释结合起来，使宏观条件预期、情绪风险感知、资金行为和流动性状态能够在统一框架中衔接。其三，在中国市场语境下，本文将宏观经济不确定性、北向资金、隐含波动率与非流动性等因素同时纳入研究，能够更系统地刻画A股核心指数收益相对宏观基本面的偏离特征，并为理解中国股票市场中"基本面—预期—偏离"的动态关系提供经验证据。
